\chapter{Leistung und Speicherbereinigung}\label{kap:leistung}

\begin{merksatz}
\boxtitel{Was die Bereinigung kostet}
\begin{center}
\renewcommand{\arraystretch}{1.5}
\begin{tabular}{@{}P{3.0cm}P{2.2cm}P{2.8cm}@{}}
\toprule
\textbf{Größe} & \textbf{Heute} & \\
\midrule
Anhaltezeit & \textcolor{gruen}{$<\mathbf{0{,}5}$ \textbf{ms}} & \textbf{meist $<100$ µs}\\
Häufigkeit & lastabhängig & \\
Anteil an der Rechenzeit & $\approx5$--$25\,\%$ & \\
Vorhersagbar & \textcolor{akzent}{\textbf{nein}} & \textcolor{akzent}{\textbf{die Grenze}}\\
\bottomrule
\end{tabular}
\end{center}

\begin{center}
\begin{minipage}{0.88\textwidth}
\kursiv{Die erste Zeile ist besser, als viele glauben --- eine halbe Millisekunde ist für fast jede Anwendung unbemerkbar. Die letzte Zeile bleibt trotzdem: \enquote{meist unter hundert Mikrosekunden} ist keine Zusicherung, und eine harte Schranke verlangt eine.}
\end{minipage}
\end{center}

\textbf{Der Vergleich mit Band 21.}

\begin{center}
\renewcommand{\arraystretch}{1.45}
\begin{tabular}{@{}P{2.4cm}P{1.6cm}P{1.8cm}P{2.8cm}@{}}
\toprule
\textbf{Schranke} & \textbf{C} & \textbf{Go} & \\
\midrule
$2$ ms hart & \textcolor{gruen}{\textbf{ja}} & \textcolor{akzent}{\textbf{nein}} & \textbf{Messerfassung}\\
$50$ ms weich & \textcolor{gruen}{ja} & \textcolor{gruen}{\textbf{ja}} & \textcolor{gruen}{Übertragung}\\
$1$ s & \textcolor{gruen}{ja} & \textcolor{gruen}{ja} & \\
\bottomrule
\end{tabular}
\end{center}

\begin{center}
\kursiv{Der Unterschied ist nicht \enquote{schnell} gegen \enquote{langsam}, sondern \enquote{zugesichert} gegen \enquote{fast immer}. Für die Messerfassung entscheidet das; für die Übertragungsschicht nicht.}
\end{center}

\begin{technik}
\boxtitel{Der wirksame Hebel: weniger belegen}
\begin{center}
{\small
\renewcommand{\arraystretch}{1.45}
\begin{tabular}{@{}P{3.2cm}P{2.0cm}P{2.8cm}@{}}
\toprule
\textbf{Maßnahme} & \textbf{Wirkung} & \\
\midrule
Abschnitt vor\-belegen & \textcolor{gruen}{\textbf{stark}} & \code{make(...,0,n)}\\
Puffer wieder\-verwenden & \textcolor{gruen}{stark} & \code{sync.Pool}\\
Werte statt Zeiger & \textcolor{gruen}{\textbf{stark}} & \textbf{kleine Strukturen}\\
Zeichenketten bauen & \textcolor{gruen}{mittel} & \code{strings.} \code{Builder}\\
Bereinigungs\-anteil stellen & \textcolor{mittel}{Speicher gegen Zeit} & \code{GOGC}\\
\bottomrule
\end{tabular}}
\end{center}

\begin{center}
\begin{minipage}{0.86\textwidth}
\kursiv{Die Reihenfolge ist nach Wirkung geordnet, und die ersten drei sind Entwurfsentscheidungen, keine Einstellungen. Wer nichts belegt, erzeugt keine Arbeit für die Bereinigung --- das ist wirksamer als jede Stellschraube.}
\end{minipage}
\end{center}
\end{technik}
\end{merksatz}

\begin{merksatz}
\boxtitel{Wo die Werte liegen}
\begin{center}
\renewcommand{\arraystretch}{1.5}
\begin{tabular}{@{}P{2.6cm}P{2.6cm}P{3.0cm}@{}}
\toprule
& \textbf{Keller} & \textbf{Halde}\\
\midrule
Wer entscheidet & \textbf{der Übersetzer} & \textbf{der Übersetzer}\\
Kosten & \textcolor{gruen}{$\approx0$} & \textcolor{akzent}{Belegung + Bereinigung}\\
Beeinflussbar & \textcolor{mittel}{mittelbar} & \\
Sichtbar & \textcolor{gruen}{\code{-gcflags=-m}} & \\
\bottomrule
\end{tabular}
\end{center}

\begin{center}
\begin{minipage}{0.88\textwidth}
\kursiv{Anders als in C entscheidet nicht der Programmierer, wo ein Wert liegt, sondern eine Analyse im Übersetzer: Was die Funktion nicht überlebt, bleibt auf dem Keller. Man kann es nicht anordnen --- aber nachsehen und den Code so ändern, dass die Analyse zum gewünschten Ergebnis kommt.}
\end{minipage}
\end{center}

\textbf{Was einen Wert auf die Halde zwingt.}

\begin{center}
\renewcommand{\arraystretch}{1.45}
\begin{tabular}{@{}P{3.2cm}P{2.2cm}P{2.6cm}@{}}
\toprule
\textbf{Ursache} & \textbf{Vermeidbar?} & \\
\midrule
Zeiger wird zurückgegeben & \textcolor{akzent}{selten} & \\
Wert in \code{interface\{\}} & \textcolor{gruen}{\textbf{oft}} & \textbf{häufigste Ursache}\\
Abschnitt wächst & \textcolor{gruen}{ja} & \textcolor{gruen}{vorbelegen}\\
Zu groß für den Keller & \textcolor{akzent}{nein} & \\
An Goroutine übergeben & \textcolor{mittel}{teils} & \\
\bottomrule
\end{tabular}
\end{center}

\begin{center}
\kursiv{Die zweite Zeile überrascht: Jeder Wert, der als leere Schnittstelle übergeben wird --- etwa an eine Ausgabefunktion --- landet auf der Halde. Eine Protokollzeile mit fünf Werten erzeugt fünf Belegungen, auch wenn sie nie ausgegeben wird.}
\end{center}

\begin{warnung}
\boxtitelr{Protokollierung als Leistungsproblem}
\begin{center}
\renewcommand{\arraystretch}{1.45}
\begin{tabular}{@{}P{2.8cm}P{1.8cm}P{1.8cm}P{2.4cm}@{}}
\toprule
\textbf{Form} & \textbf{ns/Op} & \textbf{Byte/Op} & \textbf{Abge\-schaltet}\\
\midrule
Formatiert & $\mathbf{840}$ & $\mathbf{264}$ & \textcolor{akzent}{$\mathbf{310}$ \textbf{ns}}\\
Strukturiert & $190$ & $48$ & \textcolor{gruen}{$\mathbf{9}$ ns}\\
Mit Stufenprüfung & $190$ & $48$ & \textcolor{gruen}{$\mathbf{2}$ ns}\\
\bottomrule
\end{tabular}
\end{center}

\begin{center}
\begin{minipage}{0.88\textwidth}
\kursiv{Die letzte Spalte ist die entscheidende: Was kostet eine Protokollzeile, die wegen der eingestellten Stufe gar nicht ausgegeben wird? Bei formatierter Ausgabe dreihundertzehn Nanosekunden --- die Argumente werden ausgewertet und verpackt, bevor jemand fragt, ob sie gebraucht werden.}
\end{minipage}
\end{center}

\begin{center}
\kursiv{Bei fünfhundert Messwerten je Sekunde und drei Protokollzeilen je Wert sind das eine halbe Millisekunde je Sekunde für nichts. Strukturierte Protokollierung mit Stufenprüfung senkt das auf drei Mikrosekunden.}
\end{center}
\end{warnung}
\end{merksatz}

\begin{aufgabe}
\begin{enumerate}[label=(\alph*)]
\item Senken Sie die Belegungen.
\item Bestimmen Sie die Wirkung von \code{GOGC}.
\item Untersuchen Sie, wann Go zu langsam ist.
\end{enumerate}

\loesung
\textbf{(a) Belegungen senken.}

\textbf{Die Aufgabe.} Bündeln von $500$ Messwerten zu einem Paket, $1000$ mal je Sekunde.

\begin{center}
{\small
\renewcommand{\arraystretch}{1.45}
\begin{tabular}{@{}P{2.8cm}P{1.8cm}P{2.0cm}P{2.0cm}@{}}
\toprule
\textbf{Fassung} & \textbf{µs/Op} & \textbf{Belegungen} & \textbf{Bereinigung}\\
\midrule
Naiv & $\mathbf{184}$ & $\mathbf{1503}$ & $\mathbf{22\,\%}$\\
Abschnitt vorbelegt & $61$ & $3$ & $9\,\%$\\
Puffer wiederverwendet & $\mathbf{38}$ & $\mathbf{0}$ & \textcolor{gruen}{$\mathbf{4\,\%}$}\\
\bottomrule
\end{tabular}}
\end{center}

\begin{center}
\begin{minipage}{0.88\textwidth}
\kursiv{Faktor fünf bei der Zeit --- und die letzte Spalte zeigt, wo er herkommt: Der Anteil der Speicherbereinigung fällt von zweiundzwanzig auf vier Prozent. Man hat nicht die Rechnung beschleunigt, sondern die Arbeit für die Bereinigung beseitigt.}
\end{minipage}
\end{center}

\textbf{Der wiederverwendete Puffer.}

\begin{center}
\begin{minipage}{0.9\textwidth}
\code{var pakete = sync.Pool\{}\\
\code{\ \ \ \ New: func() any \{ return make([]Messwert, 0, 500) \},}\\
\code{\}}\\
\\
\code{p := pakete.Get().([]Messwert)[:0]}\\
\code{defer pakete.Put(p)}
\end{minipage}
\end{center}

\begin{center}
\renewcommand{\arraystretch}{1.45}
\begin{tabular}{@{}P{3.0cm}P{2.2cm}P{2.8cm}@{}}
\toprule
\textbf{Eigenschaft} & \textbf{Bewertung} & \\
\midrule
Wirkung & \textcolor{gruen}{\textbf{stark}} & bei hoher Frequenz\\
Lesbarkeit & \textcolor{akzent}{schlechter} & \textcolor{akzent}{Typum\-wandlung}\\
Fehlerquelle & \textcolor{akzent}{\textbf{Puffer weitergegeben}} & \textcolor{akzent}{\textbf{Kap.~\ref{kap:nullwert}}}\\
Wann einsetzen & \textcolor{mittel}{\textbf{nach der Messung}} & \\
\bottomrule
\end{tabular}
\end{center}

\begin{center}
\kursiv{Die dritte Zeile ist die Gefahr: Wer einen zurückgegebenen Puffer noch irgendwo hält, liest später fremde Daten. Das ist genau der Fehler aus Band 21, Kapitel~3 --- die Speicherbereinigung schützt nicht davor, weil der Speicher gültig ist.}
\end{center}

\medskip
\textbf{(b) Die Wirkung von \code{GOGC}.}

\begin{center}
\renewcommand{\arraystretch}{1.45}
\begin{tabular}{@{}P{2.6cm}P{2.0cm}P{2.2cm}P{2.0cm}@{}}
\toprule
\code{GOGC} & \textbf{Speicher} & \textbf{Bereini\-gung} & \textbf{Durch\-satz}\\
\midrule
$50$ & \textcolor{gruen}{$\mathbf{120}$ MB} & \textcolor{akzent}{$18\,\%$} & $410$ k/s\\
$100$ & $210$ MB & $9\,\%$ & $\mathbf{495}$ k/s\\
$400$ & $\mathbf{680}$ MB & \textcolor{gruen}{$\mathbf{3\,\%}$} & $\mathbf{520}$ k/s\\
Speichergrenze & \textcolor{gruen}{\textbf{gedeckelt}} & \textcolor{gruen}{passt sich an} & \textcolor{gruen}{\textbf{empfohlen}}\\
\bottomrule
\end{tabular}
\end{center}

\begin{center}
\begin{minipage}{0.88\textwidth}
\kursiv{Eine reine Tauschgröße: mehr Speicher gegen weniger Bereinigungsarbeit. Die letzte Zeile ist die bessere Antwort --- statt eines Verhältnisses gibt man eine Speicherobergrenze an, und die Laufzeitumgebung regelt selbst.}
\end{minipage}
\end{center}

\begin{center}
\kursiv{Das ist besonders in Behältern wichtig: Ohne Obergrenze kennt die Laufzeitumgebung die Speichergrenze des Behälters nicht und wird von außen abgeschossen, statt häufiger zu bereinigen.}
\end{center}

\medskip
\textbf{(c) Wann Go zu langsam ist.}

\begin{center}
{\small
\renewcommand{\arraystretch}{1.5}
\begin{tabular}{@{}P{3.0cm}P{2.2cm}P{2.6cm}@{}}
\toprule
\textbf{Aufgabe} & \textbf{Go geeignet?} & \\
\midrule
Netzarbeit & \textcolor{gruen}{\textbf{sehr}} & \textbf{wartet ohnehin}\\
Textverarbeitung & \textcolor{gruen}{ja} & \\
Zahlenrechnen & \textcolor{mittel}{$\approx1{,}3\times$ C} & \textcolor{mittel}{meist genug}\\
Enge Vektorrechnung & \textcolor{akzent}{$\approx3\times$ C} & \textcolor{akzent}{keine Vektorbefehle}\\
\textbf{Harte Schranken} & \textcolor{akzent}{\textbf{nein}} & \textcolor{akzent}{\textbf{Bereinigung}}\\
\bottomrule
\end{tabular}}
\end{center}

\begin{center}
\begin{minipage}{0.88\textwidth}
\kursiv{Die dritte Zeile überrascht viele: Bei gewöhnlicher Rechenarbeit ist Go etwa dreißig Prozent langsamer als C --- nicht Faktor drei, wie man bei einer Sprache mit Speicherbereinigung erwarten könnte. Der Übersetzer erzeugt ordentlichen Maschinencode.}
\end{minipage}
\end{center}

\begin{center}
\kursiv{Die vierte Zeile ist die echte Schwäche: Der Übersetzer nutzt Vektorbefehle kaum, und es gibt keine bequeme Möglichkeit, sie anzufordern. Wer Signalverarbeitung schreibt, ist in C oder mit einer Bibliothek besser bedient.}
\end{center}

\begin{warnung}
\boxtitelr{Die Reihenfolge, wie in Band 21}
\begin{center}
\renewcommand{\arraystretch}{1.45}
\begin{tabular}{@{}P{1.0cm}P{3.4cm}P{2.0cm}P{2.2cm}@{}}
\toprule
& \textbf{Schritt} & \textbf{Aufwand} & \textbf{Gewinn}\\
\midrule
$1$ & Messen & \textcolor{gruen}{$1$ h} & --- \\
$2$ & \textbf{Blockierungen prüfen} & \textcolor{gruen}{$2$ h} & \textcolor{gruen}{\textbf{oft $2$--$5\times$}}\\
$3$ & Belegungen senken & $1$ d & $2$--$5\times$\\
$4$ & Verfahren verbessern & $2$ d & $2$--$100\times$\\
$5$ & Teil in C & \textcolor{akzent}{Wochen} & $1{,}3$--$3\times$\\
\bottomrule
\end{tabular}
\end{center}

\begin{center}
\begin{minipage}{0.88\textwidth}
\kursiv{Der zweite Schritt ist der Go-eigene und steht vorn, weil er bei nebenläufigen Programmen am häufigsten trägt: Ein Kanal ohne Puffer, ein einzelner Leser, eine Sperre mit Streit --- Faktor drei für zwei Stunden Arbeit.}
\end{minipage}
\end{center}

\begin{center}
\kursiv{Der letzte Schritt lohnt in Go noch seltener als in Band 21: Faktor $1{,}3$ für Wochen Arbeit, dazu die Kosten der Sprachgrenze aus Band 21, Kapitel~14. Meist ist die Antwort, den Bestandteil ganz in C zu schreiben --- oder es zu lassen.}
\end{center}
\end{warnung}
\end{aufgabe}
\vspace{-4pt}
\noindent Die Werk\-zeug\-do\-ku\-men\-ta\-tion\zit{27} erklärt die Pro\-fil\-ar\-ten; Cheney\zit{2} hat die Mess\-pra\-xis am gründ\-lichs\-ten be\-schrie\-ben; Wirth\zit{23} erinnert daran, dass die Da\-ten\-struk\-tur vor der Ein\-stel\-lung kommt.
